\documentclass[11pt]{article}
\usepackage{amsmath,amssymb,amsthm,hyperref}
\newtheorem{theorem}{Theorem}
\newtheorem{lemma}[theorem]{Lemma}
\title{Continuous Parameters in Merged Refinement Graphs}
\author{}
\date{November 23, 2026}
\begin{document}
\maketitle

\paragraph{Scope.}
This note treats finite graphs of coordinate refinements, allowing
arbitrary merging of equal states. It proves a parameter-dimension bound
for fixed graphs, then uses projection counting and spanning forests to
exclude first variations of all positive-step graphs, including arbitrary
node-dependent flips, at the requested rate under the stated boundary
charging convention. It does not prove or disprove the universal head
conjecture. Higher-order variations and other charging schemes remain
outside the argument.

\section{The graph and its linear realization space}

Fix $d\ge2$. A finite directed acyclic graph has internal vertices and
terminal vertices. Each internal vertex $v$ has one specified child
$c_i(v)$ for each $i\in\{0,\ldots,d-1\}$; repeated appearances of a
vertex are shared, not counted as separate copies. There are $L$ terminal
vertices, $R$ roots, and $c$ weakly connected components. No component
consists of an isolated unrefined vertex.

A realization assigns $p_v\in\mathbb R^d$ and a scalar $s_v$ at each
internal vertex, satisfying
\[
                      p_{c_i(v)}=p_v+s_v e_i.
 \tag{1}
\]
These equations define a real linear space $\mathcal V$. Let
$q=\dim\mathcal V$. An admissible positive-step realization additionally
has positive state coordinates and $s_v>0$. The inequalities do not
increase $q$. Only the root and terminal denominators are charged for
the telescoped reciprocal identity, giving at most $J=R+L$ heads.

\begin{theorem}[Merged-graph dimension bound]
For every such fixed graph,
\[
                         q\le L+(d-1)c.
 \tag{2}
\]
If it admits a positive-step realization, then each component has at
least $d$ terminals. Consequently
\[
                     q\le(2-1/d)L<2J.
 \tag{3}
\]
\end{theorem}
\begin{proof}
Map a realization linearly to the $L$ terminal numbers
\[
                       p_{v,0}-p_{v,1}.
 \tag{4}
\]
We calculate the kernel of this map. For a realization in the kernel,
write $g_v=p_{v,0}-p_{v,1}$. Equation (1) gives
\[
                    g_{c_0(v)}=g_v+s_v,\qquad
                    g_{c_1(v)}=g_v-s_v.
 \tag{5}
\]
Backward induction in the acyclic graph proves $g_v=0$ at every
vertex: it holds at the terminals; if it holds at both indicated
children, adding (5) gives $2g_v=0$. Equation (5) then gives
$s_v=0$ at every internal vertex. Equation (1) now says that $p$
is constant across every edge, hence constant on each weak component.
The component's translation vector has its first two coordinates equal.
Conversely, any such independent translations, with all $s_v=0$,
belong to the kernel. Its dimension is exactly $(d-1)c$.
Rank-nullity and the $L$-dimensional codomain give (2).

For the second assertion, choose one internal vertex in a component
of a positive-step realization. For each coordinate $i$, repeatedly
follow the child of color $i$ until reaching a terminal. Acyclicity
and finiteness guarantee termination. Along that path all coordinates
except $i$ stay fixed, while coordinate $i$ increases strictly.
Thus the $d$ endpoints have respective positions $p_v+t_i e_i$
with $t_i>0$, which are pairwise distinct. They are therefore distinct
terminal vertices. Summing over components gives $dc\le L$.
Substituting in (2) proves (3).
\end{proof}

The backward-induction argument and its consequence that all edge
variations vanish are kernel checked as
\texttt{HeadComplexity.two\_branch\_boundary\_rigidity} and
\texttt{HeadComplexity.merged\_refinement\_boundary\_rigidity}
in \texttt{Results/MergedRefinementRigidity.lean}. The graph dimension
and terminal-count arguments above are manual.

\section{Coordinate flips do not evade the dimension bound}

The argument extends to graphs in which every internal vertex chooses
its own orientation for each variable. Write $d=n+1$ and replace (1) by
\[
 p_{c_0(v)}=p_v+s_v e_0,\qquad
 p_{c_i(v)}=p_v+s_v d_{v,i}\quad(1\le i\le n),
 \tag{10}
\]
where $d_{v,i}$ is either $e_i$ or $e_0-e_i$. The choice is fixed as
part of the graph. Positive-step realizations require $s_v>0$; if a
reciprocal construction is intended, all resulting denominators must
also remain positive with an allowed coefficient orientation.

\begin{theorem}[Flipped refinements]
For the linear realization space of (10),
\[
 q\le L+nc.
 \tag{11}
\]
If a positive-step realization exists and each component has an
internal vertex, then $L\ge(n+1)c$ and $q<2J$.
\end{theorem}
\begin{proof}
Use the fixed covector $\ell=(2,1,\ldots,1)$ and map a realization
to the $L$ numbers $g_v=\ell(p_v)$ at its terminals. Both choices
of $d_{v,i}$ satisfy $\ell(d_{v,i})=1$, whereas $\ell(e_0)=2$.
For a realization in the kernel, backward induction uses
\[
 g_{c_0(v)}=g_v+2s_v,\qquad g_{c_1(v)}=g_v+s_v.
\]
Zero values at the two children force $s_v=g_v=0$. Every state
is then constant on its weak component. The kernel consists exactly
of the independent component translations in $\ker\ell$ and has
dimension $nc$. This proves (11) by rank-nullity.

For the terminal count, project states to their last $n$ coordinates.
Following only color-zero children reaches a terminal with the same
projection. Fix an internal vertex $v$. Apply this observation to
$v$ and separately to each $c_i(v)$ for $1\le i\le n$. The resulting
terminal projections are $\pi(p_v)$ and
$\pi(p_v)\mathbin{\pm}s_v e_i$, respectively. They are affinely
independent because $s_v>0$, so there are at least $n+1$ distinct
terminals in that component. This proves the remaining inequalities.
\end{proof}

The distinct-rate boundary induction and the existence of a terminal
with an invariant branch value are kernel checked as
\texttt{HeadComplexity.two\_rate\_boundary\_rigidity} and
\texttt{HeadComplexity.constant\_branch\_reaches\_terminal}.
The covector calculation, affine independence, and dimension argument
remain manual. No positivity of telescoping weights is asserted for
flipped refinements. The sign-count argument below depends only on
the realization dimension and denominator admissibility, and therefore
applies to fixed flipped graphs as well.

\section{The reciprocal identity and the sign-count scope}

Take $d=n+1$, $z_0=1$, $z_i=x_i\in\{0,1\}$, and define
\[
                B_v(x)=\sum_{i=0}^n p_{v,i}z_i,
                \qquad K_v(x)=p_{v,0}/B_v(x).
\]
At every internal state, the P\'olya identity gives
\[
            K_v=\sum_{i=0}^n\frac{p_{v,i}}{\sum_j p_{v,j}}K_{c_i(v)}.
 \tag{6}
\]
Iterating over a finite acyclic graph expresses each root feature as
a positive weighted sum of terminal features. All repeated terminals
merge. The common positive denominator orientation is preserved.
The dimension bound does not count unmerged occurrences.

For a fixed graph, all root and terminal affine coefficients are
linear functions of $q<2J$ coordinates on $\mathcal V$. Consider the
enlarged first-variation family
\[
 P_0(x)+\sum_{h=1}^J\left(\frac{D_h(x)}{B_h(x)}
                         +\frac{U_h(x)V_h(x)}{B_h(x)^2}\right),
 \tag{7}
\]
with arbitrary affine $P_0,D_h,U_h,V_h$, and positive denominators
from the graph. This enlargement includes every affine-numerator
first variation, even when the variation leaves the graph's own
linear realization space.

Set $k_1=(3J+1)(n+1)$ and $p=k_1+q$. Clearing the squared
denominators gives degree at most two in the numerator block and
at most $2J$ in the realization block. The two-block sign bound
already proved and sourced in \texttt{mechanism-audit.tex} therefore
gives, when $p\le N=2^n$,
\[
                 \left(\frac{32768N}{p}\right)^p
                           2^{k_1}(2J)^q.
 \tag{8}
\]
For $J\le C N/n^2$ its base-two logarithm is
\[
                 O_C(N\log n/n)+O_C(N/n)=o(N).
 \tag{9}
\]
Thus each fixed graph realizes only $2^{o(N)}$ strict sign tables
at the requested head rate. A family of $2^{o(N)}$ graph choices
has the same obstruction. A universal construction within this
first-variation scheme would need at least $2^{N-o(N)}$ graph
choices, or would have to leave the stated family in another way.

\paragraph{What remains unproved.}
The fixed-graph argument alone supplies no sufficiently small bound on
the number of realizable graph incidence patterns. The projection and
spanning-forest arguments below resolve this issue for all unflipped
positive-step graphs, including repeated state vectors. The later paired-projection argument also controls the union over
arbitrary flipped graphs. Higher-order perturbations remain outside
the argument.
Low continuous dimension for each member of an arbitrarily large
discrete family is not a sign-count bound for their union. No
universal labeling construction from those graph choices is supplied.

\section{A projection bound on distinct unflipped states}

The following additional constraint uses the combinatorial
Loomis--Whitney inequality. Its projection and entropy forms are recalled
in Balister and Bollob\'as, \emph{Projections, Entropy and Sumsets},
Sections 2--3, \url{https://arxiv.org/pdf/0711.1151}.

\begin{theorem}[Distinct-state bound]
Consider any finite unflipped refinement graph satisfying (1).
Let $S$ be the set of its distinct state vectors and $T\subseteq S$
the set of its distinct terminal state vectors. Then
\[
              |S|^{d-1}\le |T|^d. \tag{12}
\]
Thus $|S|\le L^{d/(d-1)}$. This does not require distinct vertices to
have distinct states. It bounds states, not the number of occurrences
of a state or the number of different refinement rules at that state.
\end{theorem}
\begin{proof}
Let $\pi_i$ delete coordinate $i$. Starting from any vertex and
following only its color-$i$ child reaches a terminal, since the graph
is finite and acyclic. Each such edge preserves $\pi_i$. Therefore
\[
                         \pi_i(S)=\pi_i(T)
                         \quad(0\le i<d). \tag{13}
\]
The inclusion from right to left uses $T\subseteq S$; the reverse
inclusion is the preceding path argument.

For completeness, the finite projection inequality follows directly
from entropy. Assume $S$ is nonempty and let $X$ be uniform on $S$.
The chain rule and decrease of conditional entropy under additional
conditioning give
\[
 \sum_{i=0}^{d-1}H(X_{[d]\setminus\{i\}})
 \ge \sum_{i=0}^{d-1}\sum_{j\ne i}
             H(X_j\mid X_0,\ldots,X_{j-1})
 =(d-1)H(X).
\]
Here the chain-rule expansion on the left conditions on the preceding
coordinates other than $i$, a subset of all preceding coordinates.
Since $H(X)=\log_2|S|$ and
$H(\pi_i(X))\le\log_2|\pi_i(S)|$, exponentiation yields
\[
              |S|^{d-1}\le\prod_{i=0}^{d-1}|\pi_i(S)|.
\]
Use (13) and $|\pi_i(T)|\le |T|$ to obtain (12). Empty $S$ is
immediate, since $d\ge2$.
\end{proof}

The path argument is kernel checked in
\texttt{HeadComplexity.refinement\_projection\_reaches\_terminal}.
The companion
\texttt{HeadComplexity.refinement\_card\_le\_terminal\_card\_sq}
proves the weaker square bound for a graph with injectively labeled
states: choose terminals reached in two different colors; the two
terminal states determine every coordinate of the starting state.
The entropy argument and its sharper exponent are manual.

For $d=n+1$ and $L\le C2^n/n^2$ with fixed $C$, (12) implies
$|S|\le(2+o(1))L$. If graph vertices have distinct states and each
internal state has one specified refinement rule, a labeled graph on
$M$ states has at most $2^M M^{dM}$ possible terminal flags and
child maps. Summing over $M\le\lfloor L^{d/(d-1)}\rfloor$ still
gives only the bound
\[
                         2^{O(nL\log(L+1))}.
\]
At the requested scale this is $2^{O_C(2^n)}$, not $2^{o(2^n)}$.
It is therefore insufficient to combine with (9) to exclude all
target-dependent graphs. Neither an upper sign-coverage theorem nor a
lower bound for the full head model follows from this enumeration.
The direct graph-count statement explicitly requires one rule per
distinct state. The next argument avoids counting complete graphs and
does not require that assumption.

\section{Spanning forests close the unflipped first-variation route}

\begin{theorem}[Arbitrary unflipped graphs]
Fix $C>0$ and put $N=2^n$. Consider all finite unflipped refinement
graphs in $d=n+1$ coordinates that admit realizations with positive
coordinates and strictly positive steps. Require every weak component
to contain an internal vertex. Charge every root and terminal occurrence,
with total $J\le CN/n^2$. Allow the graph, its number of internal
vertices, all its parameters, and the coefficients in (7) to depend
on the target. The union of the strict sign tables of (7), over all
these choices, has cardinality $2^{o(N)}$ as $n\to\infty$.
In particular this entire construction family is not universal at the
requested head rate.
\end{theorem}
\begin{proof}
For any one realization, identify vertices with the same state vector
only for purposes of encoding their positions. Keep every original edge
in the resulting geometric graph, allowing several outgoing refinement
rules at a state. This operation preserves geometric connectivity; it
need not preserve a single-rule graph presentation and we do not use
such a presentation. Every geometric edge joins points differing in one
coordinate. Positive steps exclude loops with equal endpoints.

Let $M$ be the number of distinct state vectors and $c'$ the number
of weak components of this geometric graph. The distinct-state theorem
gives $M\le J^{(n+1)/n}$. Each geometric component contains the image
of an original internal vertex. Its $d$ monochromatic paths have
distinct terminal endpoint vectors, as in the proof of (3). Components
have disjoint state sets, so $dc'\le |T|\le J$.

Choose a spanning tree in each geometric component and root each tree
arbitrarily. Label the $M$ states. Record each nonroot vertex's parent
and the coordinate changed on that tree edge. Starting from an arbitrary
vector at each root, introduce one real signed displacement per tree
edge. The resulting state vectors are linear functions of
\[
                  q_0=dc'+M-c'=M+(d-1)c'\le M+J
                  \tag{14}
\]
real parameters. Every original realization is included, because along
the chosen tree every edge is parallel to its recorded coordinate
axis. Allowing arbitrary signed displacements is an enlargement;
positivity and the other refinement constraints may be discarded when
counting sign tables. Head denominators are obtained by recording which
state occurs at each charged root or terminal occurrence. Repeated
occurrences and different rules at the same state cause no difficulty.

For fixed $C$ and sufficiently large $n$, $J^{1/n}\le3$, hence
$M\le3J$ and $q_0\le4J$. Parent maps and edge colors give at most
$(M+1)^M d^M$ encodings, including maps which are not forests as a
harmless overcount. Assigning $J$ ordered head occurrences to states
costs at most $M^J$. Summing over $1\le M\le3J$ bounds the number
of valid templates by
\[
           3J(3J+1)^{4J}d^{3J}
              =2^{O(J\log(J+1)+J\log(n+1))}. \tag{15}
\]
Only forests satisfying $dc'\le J$ are needed. Their templates have
the parameter bound (14); invalid maps counted in (15) are not added
as function families.

Within each template, use the unrestricted affine numerators in (7).
With $k_1=(3J+1)(n+1)$ and $p=k_1+q_0$, clearing the positive
squared denominators gives degree at most two in the numerator block
and at most $2J$ in the forest-parameter block. Apply (8), with $q_0$
in place of $q$. Its logarithm is
\[
 O\bigl(nJ\log(N/(nJ))+nJ+J\log(J+1)\bigr)
                  =O_C(N\log n/n)=o(N). \tag{16}
\]
For a uniform bound over smaller values of $J$, pad (7) with zero
numerators using an existing denominator, up to
$\lfloor CN/n^2\rfloor$. This does not add states or components.
The empty-score family contributes only the affine scores already
allowed by $P_0$ and is also subexponential. Thus the estimates can all
be made at the largest budget, where $p<N$ for sufficiently large $n$.
The logarithm of (15) is $O_C(N/n)=o(N)$, so taking the union of
all templates preserves $2^{o(N)}$ sign coverage.
\end{proof}

The hypothesis excluding isolated unrefined components is essential to
this argument: each isolated state would introduce an independently
parameterized denominator, invalidating the bound on root parameters. The theorem permits arbitrary numbers of internal vertex
occurrences, arbitrary real positive steps, and multiple refinement
rules at a repeated state. It imposes no integrality, bounded-depth,
or fixed-denominator-list restriction. It concerns first variations
only. The projection and forest encoding, rather than an unsupported
enumeration of all refinement graphs, are the new counting mechanism.

\section{Why the distinct-state argument does not cover arbitrary flips}

The following explicit admissible family prevents an unqualified extension
of (12) to node-dependent flips. Fix $n\ge1$, an integer $m\ge1$,
$a>0$, $b_i>0$, and $\delta>0$. For
$0\le k\le m$ and $z\in\{0,1\}^n$, use the state denominator
\[
 B_{k,z}(x)=a+k\delta+\sum_{i=1}^n(b_i+\delta z_i)x_i.
 \tag{17}
\]
The states with $k=m$ are terminal. At a state with $k<m$, let the
constant child be $(k+1,z)$. If $z_i=0$, let child $i$ be
$(k,z+e_i)$; its coefficient increment is $\delta e_i$.
If $z_i=1$, let child $i$ be $(k+1,z-e_i)$; its increment is
$\delta(e_0-e_i)$. These are exactly the allowed flipped refinement
rules in (10), with one common positive step $\delta$ at every node.
All state coefficients remain strictly positive.

The integer potential $2k+\sum_i z_i$ increases by two on a constant
edge and by one on every variable edge. Hence the graph is acyclic.
Every state is reachable from $(0,0)$ by unflipped variable steps
at level zero followed by constant steps, and this is its only root. There are
\[
                    L=2^n,\qquad M=(m+1)2^n,\qquad J=2^n+1.
 \tag{18}
\]
The states are distinct because their constant coordinates determine
$k$ and their variable coordinates determine $z$. Thus $M$ is
unbounded for fixed $n$ and $L$, and for $m\ge2$ it exceeds
$L^{(n+1)/n}=2^{n+1}$.

This is a counterexample to the proposed projection bound for arbitrary
flips, not to the universal head conjecture. Its terminal count is
$2^n$, above the requested asymptotic budget. It establishes neither
small-budget universality nor impossibility for arbitrary flipped
refinements. In particular, the forest count above must not be silently
extended to this class.

\section{A finite state bound for flips below the full-cube threshold}

Although (12) fails with arbitrary flips, a weaker bound controls distinct
states when the number of distinct terminal vectors is below $2^n$.
It is not sufficient for the sign-count conclusion of the forest theorem.

\begin{theorem}[Flipped state bound]
Let $S$ and $T$ be the distinct state and terminal sets of any finite
positive-step graph satisfying (10), put $M=|S|$, $L=|T|$, and
$N=2^n$, and assume $1\le L<N$. Then
\[
                        (N-L)M\le (N-1)L^2. \tag{19}
\]
In particular, if $L=O(N/n^2)$, then $M\le(1+o(1))L^2$.
\end{theorem}
\begin{proof}
At any state in $T$, stop even if that vector also occurs at an internal
vertex. At every state of $S\setminus T$, choose one of its original
outgoing refinement rules. This choice preserves all required children.
The functional $\ell(p)=2p_0+\sum_i p_i$ increases strictly along every
chosen edge: by $2s$ on the constant edge and by $s$ on each variable
edge. Thus every chosen path terminates in a finite set unless it stops
by a separately specified rule.

Following only constant children preserves the slope vector
$(p_1,\ldots,p_n)$ and reaches $T$. There are therefore at most $L$
distinct slope vectors in $S$. Let $K$ be the largest number of states
with one fixed slope vector. Then $M\le LK$.

For each $x\in\{0,1\}^n$, write $B_x(p)=p_0+\sum_i p_ix_i$.
A chosen rule has a unique exceptional input $x$ at which all its
variable selectors equal one. At every other input, at least one
variable child preserves $B_x$. Let $I_x$ be the set of nonterminal
states whose chosen rule has exceptional input $x$. These sets partition
$S\setminus T$, so $\sum_x|I_x|=M-L$.

Starting from any state, follow a $B_x$-preserving variable child
until reaching $T$ or $I_x$. Strict increase of $\ell$ and finiteness
ensure termination. Hence every value of $B_x$ on $S$ already occurs
on $T\cup I_x$, giving
\[
                       |B_x(S)|\le L+|I_x|.
\]
On a fixed slope fiber, distinct states have distinct constant
coordinates and hence distinct values of $B_x$. Consequently
$K\le L+|I_x|$ for every $x$. Averaging these inequalities gives
$NK\le NL+M-L$. Combine this with $M\le LK$:
\[
                   NM\le LNK\le NL^2+LM-L^2.
\]
Rearranging proves (19). The stated asymptotic consequence follows
because $L/N=O(n^{-2})$.
\end{proof}

The proof permits repeated vertices and different rules at one state.
It is a manual combinatorial theorem; no Lean declaration of (19)
is claimed. Unlike the linear-order unflipped bound, a quadratic bound
in $L$ does not make the colored-forest union subexponential at the
target head count. The paired-projection argument below supplies the sharper estimate
needed for that route. The family (17) shows that a finite bound solely
in $n,L$ is impossible once $L=2^n$.

\section{Paired projections close the arbitrary-flip first-variation route}

The coarse estimate (19) can be strengthened by using hyperplane
projections in two different coordinate bases. This removes the remaining
graph-choice gap for first variations with strictly positive steps.

\begin{theorem}[Paired-projection state bound]
Under the hypotheses of (19),
\[
             2M\le L^{1/n}(M+L),\qquad
             M\le\frac{L^{1+1/n}}{2-L^{1/n}}. \tag{20}
\]
Consequently, for each fixed $C>0$, if $L\le C2^n/n^2$ and $n$
is sufficiently large, then
\[
                       M\le \frac{2n}{\log n}L, \tag{21}
\]
where the logarithm in (21) is natural.
\end{theorem}
\begin{proof}
Choose one outgoing rule at each state of $S\setminus T$ and stop
at every state of $T$, as in the proof of (19). Let $\pi_0$ be the
quotient map with kernel $\mathbb R e_0$. The constant-child argument
gives $|\pi_0(S)|\le L$.

For $1\le i\le n$, let $\pi_i^+$ and $\pi_i^-$ be quotient maps
with respective kernels $\mathbb R e_i$ and
$\mathbb R(e_0-e_i)$. Let $a_i$ count the nonterminal states whose
chosen color-$i$ direction is $e_0-e_i$. Following the color-$i$
child while its direction is $e_i$ preserves $\pi_i^+$, and stops
at a terminal or one of these $a_i$ states. The strictly increasing
functional $\ell$ ensures termination. Following the opposite choice
similarly preserves $\pi_i^-$. Therefore
\[
       |\pi_i^+(S)|\le L+a_i,\qquad
       |\pi_i^-(S)|\le L+(M-L-a_i)=M-a_i. \tag{22}
\]

Apply the finite Loomis--Whitney inequality in the standard basis
$(e_0,e_1,\ldots,e_n)$ and separately in the basis
$(e_0,e_0-e_1,\ldots,e_0-e_n)$. The latter is a basis because its
last $n$ vectors have independent variable-coordinate projections.
A coordinate-deleted projection is, up to an invertible map of its
image, the quotient by the corresponding basis direction; hence its
cardinality is precisely the relevant cardinality in (22). Multiplying
the two projection inequalities gives
\[
 M^{2n}\le |\pi_0(S)|^2
             \prod_{i=1}^n|\pi_i^+(S)|\,|\pi_i^-(S)|
       \le L^2\prod_{i=1}^n(L+a_i)(M-a_i)
       \le L^2\left(\frac{M+L}{2}\right)^{2n}. \tag{23}
\]
The last inequality is the arithmetic-geometric mean inequality for
the two nonnegative factors whose sum is $M+L$. Taking $2n$-th
roots proves the first inequality in (20). Since $L<2^n$, its
coefficient $2-L^{1/n}$ is positive, giving the second inequality.

For the uniform asymptotic estimate, put $u=2\log n-\log C$.
For sufficiently large $n$, $u\ge\log n>0$ and $u/n\le1$.
The budget assumption implies $L^{1/n}\le2e^{-u/n}$. Since
$1-e^{-t}\ge t/2$ for $0\le t\le1$,
\[
 2-L^{1/n}\ge 2(1-e^{-u/n})\ge u/n\ge(\log n)/n,
 \qquad L^{1/n}\le2.
\]
Substitute these inequalities in (20) to obtain (21).
\end{proof}

\begin{theorem}[All positive-step refinement graphs]
Fix $C>0$. The enlarged first-variation family (7), taken over all
finite graphs (10) with arbitrary node-dependent selector flips and
strictly positive steps, has only $2^{o(2^n)}$ strict sign tables when
all root and terminal occurrences are charged and their total is at
most $C2^n/n^2$. Every weak component must contain an internal vertex.
The conclusion permits arbitrary target-dependent graphs, repeated
state vectors, different rules at a repeated state, and unbounded
numbers of internal occurrences.
\end{theorem}
\begin{proof}
Write $K=\lfloor C2^n/n^2\rfloor$ and work in sufficiently large
dimensions, so $K<2^n$. As before, collapse equal state vectors only
to encode the geometric graph, and choose a spanning forest. For each
component, following constant children from one internal vertex and
from its $n$ variable children reaches terminals with affinely independent
slope projections. Thus, if the geometric graph has $c'$ components,
$(n+1)c'\le L\le K$.

Every forest edge is parallel to one of the $2n+1$ directions
\[
                e_0,\ e_1,\ldots,e_n,\ e_0-e_1,\ldots,e_0-e_n.
\]
The root-vector and signed-edge-displacement encoding therefore has
$q_0=M+nc'\le M+K$ real parameters. By (21),
\[
                 M=O_C(2^n/(n\log n)),\qquad
                 q_0=O_C(2^n/(n\log n)). \tag{24}
\]
The number of parent maps, direction choices, and assignments of $K$
ordered head occurrences to states has logarithm
\[
                  O(M\log(M+1)+M\log(n+1)+K\log(M+1))
                       =O_C(2^n/\log n)=o(2^n). \tag{25}
\]
Summing over all eligible $M$ changes only the constant in this bound.
As before, invalid parent maps are only a combinatorial overcount;
the parameter bound is used only for valid forests. Smaller head counts
are included by zero-numerator padding with an existing denominator.

For each valid template, the numerator block of (7) has
$k_1=(3K+1)(n+1)=O_C(2^n/n)$ parameters and degree at most two.
The forest block has $q_0$ parameters and degree at most $2K$
after clearing squared denominators. Thus the logarithm of the
bound (8), with $p=k_1+q_0$ and $q_0$ in place of $q$, is at most
\[
        O_C(2^n\log n/n)+O(q_0\log(K+1))
                   =O_C(2^n\log n/n+2^n/\log n)=o(2^n). \tag{26}
\]
Here $p<2^n$ for sufficiently large $n$, as required by (8).
Combine (25) and (26) to obtain the asserted union bound.
\end{proof}

These are manual theorems. The kernel-checked constant-branch reachability
lemma supports the stopping arguments, but the quotient-projection,
entropy, forest, and sign-count deductions have not been formalized.
The ladder (17) is consistent with (20): its $L=2^n$ makes the
denominator in (20) zero. The exclusion of unrefined isolated components
and the restriction to first variations remain essential. Higher-order
perturbations, nonpositive steps, and arbitrary admissible head lists
are not excluded by this theorem. Proportional but unequal state vectors
are not identified in the charged boundary: for example $(1,1)$ and
$(2,2)$ give the same normalized reciprocal feature but remain distinct
vectors. A construction charging only projective classes after combining
proportional fractions requires a separate argument. No affirmative
worst-case upper bound is obtained.

\section{Finite structural probe}

The script \texttt{probe\_merged\_polya\_dimension.py} generates
finite positive-step graphs by merging equal integer states. Coordinate
equalities are first contracted, after which the remaining realization
equations have integer coefficients. Rank over the recorded prime gives
an upper bound on the real kernel dimension. Any potential excess over
$J$ is scheduled for rational rank and boundary-projection verification.

The recorded run examined $126839$ graphs and found no excess over
the stronger, unproved bound $q\le J$. That finite outcome is not
used in the proof of (2), does not establish $q\le J$, and does not
prove any unrestricted head upper or lower bound. The search was
stopped after the independent proof of (2) supplied the needed
linear-order parameter estimate for each fixed graph.

\end{document}
